\documentclass[11pt]{article}
\usepackage{amsmath,amssymb,amsthm}
\usepackage[margin=1.1in]{geometry}
\newtheorem{theorem}{Theorem}
\newtheorem{corollary}[theorem]{Corollary}
\newtheorem{remark}[theorem]{Remark}
\newcommand{\Sym}{\mathrm{Sym}}
\newcommand{\eps}{\varepsilon}
\title{The root number of the tensor tower of an elliptic curve:\\
a closed form, a Kummer criterion, and odd-order central vanishing\\
at the Fermat grades}
\author{Samuel Lavery\thanks{With machine-checked components in Lean~4
(\texttt{RootNumberLaw.lean}, \texttt{TensorTowerGlobalSign.lean}); axiom footprint
$\{\texttt{propext},\texttt{Classical.choice},\texttt{Quot.sound}\}$, no \texttt{sorry}.}}
\date{July 31, 2026 --- draft}
\begin{document}
\maketitle

\begin{abstract}
Let $E/\mathbb Q$ be an elliptic curve and $M_g=H^1(E)^{\otimes g}$ the $g$-fold tensor
motive.  We give the archimedean epsilon factor of the tower in closed form: with
$A(g)=2^{\,g-1}+g\binom{g-1}{\lfloor (g-1)/2\rfloor}$ one has
$\eps_\infty(M_g)=i^{A(g)}$, the exponent $A(g)$ is even at every grade (the sign is
real), $\eps_\infty=+1$ at every even grade, and for odd $g\ge3$
\[
\eps_\infty(M_g)=-1 \iff g=2^k+1 ,
\]
by Kummer's carry criterion for the $2$-adic valuation of a central binomial
coefficient.  For semistable $E$ we then assemble the \emph{global} root number of
$L(M_g,s)$ across the Clebsch--Gordan decomposition and prove that at every Fermat
grade $g=2^k+1$ $(k\ge2)$ the global sign is $-1$ \emph{for both parities of}
$\eps(E)$: $L(M_g,s)$ vanishes at $s=(g+1)/2$ to \emph{odd} order at least
$M(g,\eps(E))$, an explicitly computed ledger ($9$ and $1$ at $g=5$ for the two
parities; $91$ and $35$ at $g=9$; $\approx1.9\times10^{18}$ by $g=65$).  At $g=3$
the ledger recovers the Gross--Schoen grade ($M=3$ for $\eps(E)=-1$), matching the
sign of the triple-product $L$-function (Gross--Kudla); above $g=3$ the tensor
tower's sign appears not to have been computed before.  The combinatorial laws are
machine-checked (exact ledger values through $g=9$ and sign parities through
$g=25$ in the kernel); the arithmetic inputs are the classical local--global sign
computations of Deligne, Martin--Watkins, Dummigan--Martin--Watkins and Saito, and
the functional equations are unconditional by symmetric-power functoriality
(Newton--Thorne).
\end{abstract}

\section{The object and what is known}

Write $V=H^1(E)$, a motive of weight $1$ with $\det V=\mathbb Q(-1)$, and
$M_g=V^{\otimes g}$, of weight $g$ and rank $2^g$.  Its $L$-function factors through
the Clebsch--Gordan decomposition
\begin{equation}\label{eq:cg}
V^{\otimes g}\;=\;\bigoplus_{j=0}^{\lfloor g/2\rfloor}
\bigl(\Sym^{\,g-2j}V\bigr)(-j)^{\oplus m_{g,j}},\qquad
m_{g,j}=\binom gj-\binom g{j-1},
\end{equation}
the multiplicities being the ballot numbers (the decomposition holds at
Chow-motive level \cite{Cao}), so that
$L(M_g,s)=\prod_j L(\Sym^{g-2j}E,\,s-j)^{m_{g,j}}$.  Every twisted factor's completed
functional equation reflects about the same axis: the center of the piece
$\Sym^{g-2j}(-j)$ sits at $s=\frac{g-2j+1}2+j=\frac{g+1}2$, the center of the tensor
motive.  By symmetric-power functoriality for holomorphic newforms
\cite{NewtonThorne} (non-CM; the CM case is classical via Hecke characters) each
factor is automorphic, so $L(M_g,s)$ has entire continuation (odd $g$) and its
functional equation is a theorem, not a hypothesis.

The literature on the \emph{sign} of the tower stops at $g=3$.  At $g=1$ the sign
theory is classical.  At even $g$ the sign is $+1$: for $g=2$ by the
symplectic$\times$symplectic mechanism (Prasad--Ramakrishnan \cite[Prop.~2.1]{PR},
the argument due to Rohrlich \cite{RohrlichWD}; cf.\ \cite[Prop.~6.1]{PR} for
$\Sym^2\!\otimes\Sym^2$), and for even $g\ge4$ via \eqref{eq:cg} from the local
$+1$ at every even symmetric power \cite[Prop.~7.1]{DMW} with Saito's global
theorem \cite{Saito}.  At $g=3$ the global sign of the triple product is computed
by Gross--Kudla \cite{GrossKudla} for newforms of common square-free level (their
(1.9) is exactly the Hodge multiset of $V^{\otimes3}$, and the leading $-1$ of
their sign formula is the archimedean factor), with the local archimedean
dichotomy in Prasad \cite{PrasadTri}.  For $g\ge4$ we are not aware of any
computation of the sign of the tensor tower, nor of any appearance of $2$-adic
binomial criteria in the root-number literature.

\section{The archimedean law}

Deligne's theory of local constants \cite{DeligneLNM}, in the form of the
archimedean recipe \cite[\S5.3]{DeligneAnvers}, assigns to a Hodge structure its
archimedean epsilon factor through the Hodge types.  Applied to \eqref{eq:cg} --- or
directly to the Hodge diamond of $M_g$, whose type-$(p,q)$ multiplicity is
$\binom gp$ with $p+q=g$ --- the exponent of $i$ is
\[
A(g)\;=\;\sum_{p<g/2}\,(q-p+1)\binom gp\;+\;\tfrac12\binom g{g/2}\Big|_{g\ \mathrm{even}} .
\]
For odd $g$ no $(p,p)$ block occurs, so everything below is independent of the
convention for the balanced middle block.

\begin{theorem}[archimedean sign law; machine-checked]\label{thm:arch}
For all $g\ge1$:
\emph{(i)} $A(g)=2^{\,g-1}+g\binom{g-1}{\lfloor(g-1)/2\rfloor}$;
\emph{(ii)} $A(g)$ is even, so $\eps_\infty(M_g)=i^{A(g)}\in\{\pm1\}$;
\emph{(iii)} $\eps_\infty=+1$ for every even $g$; and for odd $g\ge3$,
\[
\eps_\infty(M_g)=-1\iff A(g)\equiv2\ (\mathrm{mod}\ 4)\iff g=2^k+1,\ k\ge1 .
\]
\end{theorem}

\begin{proof}[Proof (sketch; full proof in Lean)]
(i) The deviation part $\sum_p (g-2p)\binom gp$ telescopes to
$g\binom{g-1}{\lfloor(g-1)/2\rfloor}$ and the flat part is a half row, $2^{g-1}$.
(ii) is the parity of a central binomial.  (iii) For odd $g=2m+1$ the condition
$A\equiv2\ (4)$ reduces to $v_2\binom{2m}m=1$, and by Legendre's formula
$v_2\binom{2m}m=s_2(m)$, the binary digit sum; $s_2(m)=1$ iff $m$ is a power of
two, i.e.\ $g=2^k+1$.  This is Kummer's carry criterion: the sign flips exactly
when adding $m+m$ in base $2$ carries once.  The Lean formalisation
(\texttt{RootNumberLaw.lean}) proves (i)--(iii) unconditionally as statements about
$A$; the identification of $A$ with Deligne's exponent is the cited classical
recipe.  As an independent check, assembling the per-piece archimedean signs
$\eps_\infty(\Sym^m)=-\bigl(\tfrac{-2}m\bigr)$ \cite[\S5.3]{DeligneAnvers} across
\eqref{eq:cg} reproduces $i^{A(g)}$ at $g=3,5,7$ and the closed form through
$g=11$.
\end{proof}

The even half of Theorem~\ref{thm:arch} recovers what the Stiefel--Whitney mechanism
knows \cite{PR,RohrlichWD}; the odd half is, to our knowledge, new above $g=3$, where
it recovers Gross--Kudla's $-1$.

\section{The global law for semistable curves}

For semistable $E/\mathbb Q$ and odd $m$ the global sign of $L(\Sym^mE,s)$ is
\begin{equation}\label{eq:symsign}
w_m=\Bigl(\tfrac{-2}m\Bigr)\,\eps(E),
\end{equation}
assembled from Deligne's archimedean factor and $w_m(p)=w_1(p)^m$ at multiplicative
$p$ \cite[\S4.3]{MartinWatkins}, \cite{DMW}; even powers have sign $+1$
\cite{Saito}.  The sign of a factor is unchanged by the shift to the common
center (the Hodge gaps $q-p$ are unchanged, and the unramified twist scales the
finite epsilon factors by a positive real), so the global sign of the tensor
tower is the ballot-weighted product of
\eqref{eq:symsign}, and each factor with $w=-1$ vanishes to odd order at the common
center.  Define the \emph{forced multiplicity ledger}
\[
M(g,\eps_E)\;=\;\sum_{j\,:\,w_{g-2j}=-1} m_{g,j},
\qquad\text{so}\qquad
\mathrm{ord}_{s=\frac{g+1}2}\,L(M_g,s)\;\ge\;M(g,\eps_E),
\]
and the global sign is $(-1)^{M(g,\eps_E)}$.

\begin{theorem}[Fermat grades force vanishing for every semistable curve]\label{thm:global}
Let $E/\mathbb Q$ be semistable and $g=2^k+1$ with $k\ge2$.  Then $M(g,+1)$ and
$M(g,-1)$ are both odd: the global root number of $L(M_g,s)$ is $-1$ regardless of
$\eps(E)$, and $L(M_g,s)$ vanishes at $s=(g+1)/2$ to order at least
$M(g,\eps(E))\ge1$.
\end{theorem}

\begin{proof}
Write $S_-=\sum_{(\frac{-2}{g-2j})=-1}m_{g,j}$ and $S_+$ for the complementary sum,
so $M(g,+1)=S_-$, $M(g,-1)=S_+$, and by the telescoping of the ballot numbers
$S_++S_-=\binom g{(g-1)/2}$.  For $k\ge3$ we have $g\equiv1\ (\mathrm{mod}\ 8)$, and
as $j$ increases the class of $g-2j$ mod $8$ cycles $(1,7,5,3)$, so the
$(\frac{-2}\cdot)=-1$ residues occupy the consecutive pairs $j\equiv1,2\
(\mathrm{mod}\ 4)$.  Within each pair the ballot numbers telescope:
$m_{g,4t+1}+m_{g,4t+2}=\binom g{4t+2}-\binom g{4t}$, and since
$(g-1)/2=2^{k-1}\equiv0\ (\mathrm{mod}\ 4)$ the pairs are complete.  Hence
\[
S_-\;\equiv\;\sum_{\substack{0\le j\le 2^{k-1}-2\\ j\ \mathrm{even}}}\binom gj
\pmod 2 .
\]
By Lucas' theorem at $p=2$, $\binom{2^k+1}j$ is odd iff the binary digits of $j$
are dominated by $10\cdots01_2$, i.e.\ $j\in\{0,1,2^k,2^k+1\}$; in the range only
$j=0$ contributes, so $S_-\equiv1$.  The total $\binom g{2^{k-1}}$ is even by
Kummer (adding $2^{k-1}$ and $2^{k-1}+1$ carries at bit $k-1$), so
$S_+=\binom g{2^{k-1}}-S_-$ is odd as well.  The case $k=2$ ($g=5$) is a direct
kernel check.  Each sign-$(-1)$ factor $L(\Sym^{g-2j}E,s-j)$ vanishes to odd order
$\ge1$ at the common center and occurs $m_{g,j}$ times.
\end{proof}

The sign behaviour is sharp.  Central vanishing itself is generic here:
$M(g,\eps_E)>0$ for every odd $g\ge5$ and both parities.  What distinguishes the
Fermat grades is \emph{odd-order} vanishing --- global sign $-1$ --- for both
parities at once (sign parities kernel-checked through $g\le25$; the grades
$g=33,65$ are covered by Theorem~\ref{thm:global}):

\begin{center}
\begin{tabular}{c|c c|l}
$g$ & $M(g,-1)$ & $M(g,+1)$ & global sign $-1$ for\\ \hline
$3$ & $3$ & $0$ & $\eps(E)=-1$ only (Gross--Schoen grade)\\
$5$ & $9$ & $1$ & both (Fermat) --- incl.\ curves with no\\
& & & \quad sign-forced vanishing below $\Sym^5$\\
$7$ & $28$ & $7$ & $\eps(E)=+1$ only\\
$9$ & $91$ & $35$ & both (Fermat)\\
$13$ & $1078$ & $638$ & neither\\
$17$ & $14179$ & $10131$ & both (Fermat)\\
$25$ & $2848366$ & $2351934$ & neither\\
$33$ & $623067171$ & $543735939$ & both (Fermat)\\
$65$ & $1863283301142158627$ & $1746430915865974243$ & both (Fermat)
\end{tabular}
\end{center}

\begin{corollary}
For every semistable $E/\mathbb Q$ and every $k\ge2$, the tensor-tower
$L$-function $L(H^1(E)^{\otimes(2^k+1)},s)$ vanishes at its center to \emph{odd}
order: the global root number is $-1$ regardless of $\eps(E)$.  Since central
vanishing of even order occurs generically at odd grades, the Fermat grades are
exactly the displayed grades where the sign forces vanishing for every semistable
curve at once.  Under the Beilinson--Bloch philosophy each forced odd-order zero
predicts null-homologous cycles on $E^{2^k+1}$ generalising the Gross--Schoen
modified diagonal (the $g=3$ row, where the forcing occurs for $\eps(E)=-1$ and
$M=3$ matches the classical triple-product case \cite{GrossKudla}).
\end{corollary}

\begin{remark}[what is new, what is cited]
The closed form of Theorem~\ref{thm:arch}, the Kummer characterisation, the ballot
assembly of Theorem~\ref{thm:global} and its Fermat-parity law appear to be new; the
per-piece signs \eqref{eq:symsign}, Deligne's recipe, Saito's even-power theorem and
the Newton--Thorne functional equations are the cited inputs.  The combinatorial
statements are machine-checked in Lean~4 at standard axiom footprint: the closed
form, the parity and digit laws, the ballot telescoping, the exact ledger values at
$g=3,5,9$, and the sign parities at $g=7,13,17,25$ are kernel (\texttt{decide})
theorems.  The parities at $g=33,65$ require no computation --- they are Fermat
grades, covered by Theorem~\ref{thm:global} --- and the displayed magnitudes at
$g=13,25,17,33,65$ are evaluator outputs of the same compiled definition (the
$g=13$ values independently hand-verified in audit).  The forced multiplicities are lower bounds from
sign parity alone; no part of this note assumes any conjecture.
\end{remark}

\begin{thebibliography}{9}
\bibitem{Cao} J. Cao, \emph{Motives for an elliptic curve}, arXiv:1708.08092.
\bibitem{DeligneLNM} P. Deligne, \emph{Les constantes des \'equations
fonctionnelles des fonctions $L$}, LNM 349 (1973), 501--597.
\bibitem{DeligneAnvers} P. Deligne, \emph{Valeurs de fonctions $L$ et p\'eriodes
d'int\'egrales}, PSPM 33.2 (1979), 313--346.
\bibitem{DMW} N. Dummigan, P. Martin, M. Watkins, \emph{Euler factors and local root
numbers for symmetric powers of elliptic curves}, Pure Appl. Math. Q. 5 (2009),
1311--1341.
\bibitem{GrossKudla} B. Gross, S. Kudla, \emph{Heights and the central critical
values of triple product $L$-functions}, Compositio Math. 81 (1992), 143--209.
\bibitem{MartinWatkins} P. Martin, M. Watkins, \emph{Symmetric powers of elliptic
curve $L$-functions}, ANTS-VII, LNCS 4076 (2006), 377--392.
\bibitem{NewtonThorne} J. Newton, J. Thorne, \emph{Symmetric power functoriality
for holomorphic modular forms}, Publ. Math. IH\'ES 134 (2021), 1--116, 117--152.
\bibitem{PrasadTri} D. Prasad, \emph{Trilinear forms for representations of
$\mathrm{GL}(2)$ and local $\eps$-factors}, Compositio Math. 75 (1990), 1--46.
\bibitem{PR} D. Prasad, D. Ramakrishnan, \emph{On the global root numbers of
$\mathrm{GL}(n)\times\mathrm{GL}(m)$}, PSPM 66.2 (1999), 311--330.
\bibitem{RohrlichWD} D. Rohrlich, \emph{Galois theory, elliptic curves, and root
numbers}, Compositio Math. 100 (1996), 311--349.
\bibitem{Saito} T. Saito, \emph{The sign of the functional equation of the
$L$-function of an orthogonal motive}, Invent. Math. 120 (1995), 119--142.
\end{thebibliography}
\end{document}
